%==============================================================================
% Section 7: Discussion and Outlook
%==============================================================================
\section{Discussion and Outlook}
\label{sec:discussion}

This section organises the implications of the results and the directions
for future work.

\subsection{What This Paper Has Fixed}
\label{sec:discussion:fixed}

This paper fixes a fairly strong dictionary---rather than the bare
minimum---for placing the four Thurston geometries within a single
comparison framework.  In the first layer the 36-entry bulk inventory closes,
making explicit that geometry-dependent activation rules ride on top of
Maxwell universality.  In the second layer the three observable families
spectral / torsional-charge / inflow are clearly separated, and the pairwise
cobordism dictionary organises all six pairs into a comparable form.

This organisation positions $\Nilthree$ as the geometry that carries a
strict local APS spectral core, $\Sthree$ as the geometry that carries a
frame-bundle-normalised torsional-charge core, $\Tthree$ as the trivial
spectral benchmark, and $\Solthree$ as the geometry that may carry a global
spectral branch depending on the compact quotient and spin structure.  An
important point is that within the spectral family one must still
distinguish sources: the value in $\Nilthree$ is fixed by the local
Levi--Civita spinor CS, while the spectral entry in $\Solthree$ is raised
by global kernel / spin data.  The scaffold layer is needed to describe
this difference of source together with the distinctness of the background
geometries simultaneously.

By aligning a KK reading and a Matsubara-style thermal reading on the same
Euclidean circle we have also made explicit that, with the main dictionary
preserved, a reduced description and a thermal description can be added.
This shows that the principal result of this paper is the Euclidean
response dictionary itself, while the two readings are secondary
interpretations of it.

\subsection{Future Directions}
\label{sec:discussion:future}

Direct extensions split into at least the following three lines.

\begin{enumerate}
  \item \textbf{Lorentzian extension.}  Lifting the Euclidean response
        dictionary fixed here to a Lorentzian setting may make it possible
        to decide which protected / activated structure survives as a
        real-time observable.  In particular, which of parity-odd transport,
        interface response, and thermal interpretation will withstand
        analytic continuation is the next-stage question to be made
        explicit.

  \item \textbf{Application to phenomenology.}  The dictionary is not
        confined to abstract geometry classification but can serve as a
        platform that organises which response channels are expected for
        phenomenological labels such as metastable, rolling, or no-support
        branches.  Using bulk activation and boundary observables together
        leaves room for organising geometry dependence that is invisible
        from a single effective coefficient.

  \item \textbf{Generalisation of $\Solthree$ compact-quotient / spin-structure
        dependence.}  This paper makes the boundary spectral entry of
        $\Solthree$ strict on a compact mapping-torus benchmark, but how
        the kernel structure changes for general hyperbolic monodromies and
        quotient families is left for future work.  Making explicit how the
        local/global dichotomy fixed here generalises across the Sol family
        is the natural next step for handling borderline spectral entries.
\end{enumerate}
